\begin{textsample}{NEG Dim 1 – 2009\_c6 – Score -3.00 – 2009\_c6/enade\_2009\_publi\_15f3034279.txt}  \label{ex:f1_neg_216}
Capítulo 3 Análise Técnica da Prova Este capítulo tem por objetivo apresentar o desempenho dos estudantes de COMUNICAÇÃO SOCIAL – PUBLICIDADE E PROPAGANDA no ENADE/2009.
Para isso, foram calculadas as estatísticas básicas da prova como um todo, bem como as estatísticas das questões relacionadas à Formação Geral e ao Componente Específico.
Nas Tabelas, são apresentadas as seguintes estatísticas : total da população de presentes, média, erro-padrão da média, desvio-padrão, mediana, nota máxima, nota mínima e coeficiente de assimetria ( ver definição em 1.4.3 ).
As estatísticas apresentadas neste capítulo contemplam, separadamente, os ingressantes, os concluintes e o total de estudantes, e foram calculadas tendo -se em vista as seguintes agregações : ( a ) Região Geográfica e País e ( b ) a Categoria Administrativa e a Organização Acadêmica.
As notas brutas das questões discursivas em \textbf{cada} uma das partes, Formação Geral/Componente Específico, foram calculadas a partir da média simples das notas de \textbf{cada} uma das questões correspondentes.
A primeira questão discursiva da Formação Geral foi desconsiderada e, portanto, a nota bruta referente às questões discursiva de Formação Geral foi igual à nota da questão mantida, a segunda.
No Componente Específico, a nota bruta correspondeu à média das notas nas três questões.
As notas brutas das questões objetivas de \textbf{cada} uma das partes, Formação Geral/Componente Específico, foram calculadas a partir dos acertos dos estudantes considerando -se uma regra de três para a conversão para a escala de 0 a 100, descontando -se aquelas questões que foram desconsideradas pela comissão ou que não passaram pelo teste do índice de discriminação ponto-bisserial ( ver seção 1.4.6 para a descrição do índice ).
Questões com índices de discriminação inferior a 0,2 foram excluídas do cômputo da nota do estudante.
As notas brutas de Formação Geral e no Componente Específico para \textbf{cada} um dos estudantes foram calculadas como médias ponderadas das suas notas nas questões objetivas e discursivas, ponderação diferente para \textbf{cada} parte : C\_FG = C\_FG^OBJ × 0,6 + C\_FG^DIS × 0,4 C\_CE = C\_CE^OBJ × 0,85 + C\_CE^DIS × 0,15 onde C\_FG é a nota bruta da Formação Geral, C\_FG^OBJ é a nota bruta nas questões objetivas da Formação Geral, C\_FG^DIS é a nota bruta nas questões discursivas da Formação Geral, C\_CE é a nota bruta do Componente Específico, C\_CE^OBJ é nota bruta nas questões objetivas do Componente Específico, e C\_CE^DIS é nota bruta nas questões discursiva do Componente Específico.
A nota bruta da prova, C\_GER, por sua vez, é também uma média ponderada da nota de \textbf{cada} uma das partes, Formação Geral/Componente Específico : C\_GER = C\_FG × 0,25 + C\_CE × 0,75. 3.1 Estatísticas Básicas da Prova 3.1.1 Estatísticas Básicas Gerais A Tabela 3.1 apresenta as estatísticas básicas da prova por grupo de estudantes.
O total de estudantes inscritos para a prova foi de 36585.
Destes, 74,9 \% compareceram, sendo que a presença foi maior entre concluintes ( 85,2 \% ) do que entre ingressantes ( 66,4 \% ).
A média geral da prova foi de 44,02, sendo que os ingressantes obtiveram média mais baixa, de 40,87, que os concluintes, cuja média foi 46,98.
O desvio-padrão geral foi de 16,3, sendo o do grupo de ingressantes menor ( 15,2 ), indicando que, em contraposição ao grupo dos concluintes ( 16,7 ), aqueles possuíram uma distribuição mais homogênea nas notas.
A nota máxima foi 92,4, obtida por pelo menos um concluinte ao passo que a maior nota obtida por um ingressante foi 88,7.
Os coeficientes de assimetria para os dois grupos considerados, o de concluintes e o de ingressantes, bem como para a população como um todo foram todos negativos, indicando uma maior concentração à esquerda da distribuição e um maior espalhamento à direita.
Uma comparação entre o desempenho de ingressantes e concluintes pode ser realizada a partir dos dados apresentados no Gráfico 3.1.
Cumpre destacar que as notas dos concluintes, de um modo geral, foram mais elevadas do que as dos ingressantes, cuja distribuição de notas \textbf{está} à esquerda dos primeiros.
A ideia de melhor desempenho é reforçada pelas notas medianas dos dois grupos : 48,0 entre os concluintes e 41,2 entre os ingressantes.
Nas faixas de notas até 50 pontos o percentual de ingressantes é maior que o dos concluintes, situação esta que se inverte para as faixas de notas superiores a 50 pontos.
Constatou -se que, aproximadamente 28,3 \% dos ingressantes e 45,4 \% dos concluintes situam -se nas faixas de notas superiores a 50,0 pontos O Gráfico 3.2 apresenta informações referentes ao desempenho geral de ingressantes e concluintes, comparando os resultados considerando a Região, Categoria Administrativa e Organização Acadêmica.
Levando -se em conta as notas médias dos estudantes em \textbf{cada} Região, observa -se que, há nos dois grupos, uma variação expressiva entre as Regiões.
As notas médias mais elevadas entre os concluintes ocorreram nas Regiões Sul ( 54,03 ) e Nordeste ( 47,32 ), ambas acima da média nacional ( 46,98 ).
Em relação aos ingressantes, as médias mais altas foram obtidas também nas Regiões Sul ( 44,79 ) e Nordeste ( 41,31 ), nas demais Regiões a média foi inferior à média nacional ( 40,87 ).
Na Região Norte, foi obtida a média de notas mais baixa entre os concluintes ( 38,63 ) e ingressantes ( 36,59 ).
Cursos nesta área só foram oferecidos em instituições federais, estaduais, municipais e privadas.
Levando -se em conta os agrupamentos dos estudantes em categorias administrativas, observa -se que a pontuação média mais elevada entre os concluintes e ingressantes foi encontrada em instituições de origem estadual ( respectivamente 54,08 e 54,31 ) – ambas acima da média nacional.
Tendo como foco as notas médias dos estudantes agrupados por Organização Acadêmica, notamos que Universidades apresentaram as maiores médias entre os concluintes ( 47,46 ) e Institutos Federais de Educação, Ciência. e Tecnologia apresentaram as maiores médias entre os ingressantes ( 47,81 ).
Ao passo que as piores notas médias foram observadas nos Centros Universitários para concluintes ( 45,78 ) e ingressantes ( 39,80 ). 3.1.2 Estatísticas Básicas de Formação Geral A Tabela 3.2 apresenta as estatísticas básicas de Formação Geral por grupo de estudantes.
O total de estudantes inscritos para a prova foi de 36585.
Destes, 74,9 \% compareceram, sendo que a presença foi maior entre concluintes ( 85,2 \% ) do que entre ingressantes ( 66,4 \% ).
A média geral da parte da prova de Formação Geral foi de 39,94, sendo que os ingressantes obtiveram média mais baixa, 38,47, que os concluintes, cuja média foi 41,32.
O desvio-padrão geral foi de 20,3, sendo o do grupo de ingressantes menor ( 20,0 ), indicando que, em contraposição ao grupo dos concluintes ( 20,4 ), aqueles possuíram uma distribuição mais homogênea nas notas.
A nota máxima foi 91,4, obtida por estudantes de ambos os grupos.
Para esta componente, assim como para a nota geral, os coeficientes de assimetria foram positivos para todos os grupos.
O Gráfico 3.3 permite a comparação entre concluintes e ingressantes relativamente ao desempenho na parte da prova que aborda a Formação Geral.
Nesta parte, se excluídas as provas com nota zero, tal como na prova como um todo, também os concluintes apresentaram desempenho superior, embora reste perceptível um maior equilíbrio.
A mediana dos concluintes foi superior à dos ingressantes, o mesmo acontecendo com a média.
Dos ingressantes, aproximadamente 28,5 \% alcançaram nota superior a 50 ; dentre os concluintes, o percentual dos estudantes nesta faixa de notas foi igual a 33,3 \%.
No Gráfico 3.4 são apresentadas as informações referentes ao desempenho dos ingressantes e dos concluintes na parte de Formação Geral, em diferentes agregações como Região, Categoria Administrativa e Organização Acadêmica.
Considerando -se as notas médias dos estudantes segundo \textbf{cada} Região, observou -se que, tanto no grupo dos ingressantes quanto no grupo dos concluintes, houve diferenças marcantes.
As médias variam, entre os concluintes, de 30,99 na Região Norte a 49,31 na Região Sul.
Entre os ingressantes, os valores foram de 31,48 na Região Norte a 43,75 na Região Sul.
A maior dispersão de notas dos estudantes de uma dada Região foi encontrada no Norte tanto para os concluintes quanto para os ingressantes.
Nas Regiões Nordeste, Sudeste, Sul e Centro-Oeste os concluintes apresentaram médias maiores do que os ingressantes, a situação inversa aconteceu no Norte.
Todas as Regiões apresentaram concluintes.
Tendo como foco as categorias administrativas, destacou -se que todas as instituições \textbf{estavam} representadas.
Sendo que as notas médias das estaduais foram superiores para ingressantes ( 61,99 ) e concluintes ( 49,83 ) ; As notas mais baixas para ambos os grupos foram encontradas nas instituições privadas.
Os estudantes das instituições privadas se apresentam mais homogêneos, já que o desvio padrão para concluintes e ingressantes é o menor dos valores encontrados.
Considerando -se o tipo de Organização Acadêmica, as Faculdades apresentaram o maior valor para concluintes ( 41,87 ).
Os Institutos Federais de Educação, Ciência e Tecnologia apresentaram o maior valor para ingressantes ( 40,93 ), sendo que estes não apresentaram concluintes.
Os estudantes das Universidades se apresentaram mais homogêneos.
Os Centros Federais de Educação Tecnológica não \textbf{estavam} representados nesta área. 3.1.3 Estatísticas Básicas do Componente Específico A Tabela 3.3 apresenta as estatísticas básicas referentes ao Componente Específico da prova, tendo sua média geral, de 45,44, superior ao desempenho da prova de COMUNICAÇÃO SOCIAL – PUBLICIDADE E PROPAGANDA como um todo, cuja média foi 44,02, como apresentado na Tabela 3.1, que por sua vez foi superior ao desempenho na Formação Geral, que teve média igual a 39,94, conforme mostrado na Tabela 3.2.
Os concluintes apresentaram um desempenho médio ( 48,92 ) superior ao dos ingressantes ( 41,73 ), havendo ainda indícios de maior homogeneidade entre as notas médias dos ingressantes do que entre as dos concluintes, como sugeriu o desvio padrão, menor no grupo dos ingressantes do que no dos concluintes.
As notas máximas obtidas por concluintes ( 93,8 ) foram superiores às obtidas por ingressantes ( 92,0 ).
As distribuições de notas para \textbf{cada} grupo se apresentaram com uma inclinação para notas maiores, como evidenciados pelos coeficientes de assimetria negativos.
A situação de assimetria é confirmada pelo Gráfico 3.5 que proporciona uma comparação do desempenho entre concluintes e ingressantes em relação ao Componente Específico da prova.
Dentre as três comparações já feitas ( confrontar com Gráficos 3.1 e 3.3 ) observa -se que, neste componente, a diferença de desempenho entre concluintes e ingressantes é positivamente a favor dos concluintes.
Claramente a distribuição das notas médias dos ingressantes, à esquerda daquela dos concluintes, evidenciou o percentual mais elevado de notas médias nas faixas mais baixas para este grupo.
As medianas das distribuições das notas dos ingressantes e dos concluintes foram 42,3 e 50,3, respectivamente, enfatizando o melhor desempenho dos concluintes nesta parte da prova.
Entre os ingressantes, 31,7 \% dos estudantes obtiveram notas superiores a 50 e entre os concluintes, 50,9 \%.
Assim como os Gráficos 3.2 e 3.4, o Gráfico 3.6 apresenta uma comparação dos resultados em relação às Regiões, Categoria Administrativa e Organização Acadêmica, agora levando em conta o desempenho dos ingressantes e dos concluintes no Componente Específico da prova.
Tendo como foco as notas médias dos estudantes, em determinada Região, observa -se que, entre os concluintes, as mais elevadas foram encontradas nas Regiões Nordeste ( 49,66 ) e Sul ( 55,65 ), valores acima da média nacional para esta componente ( 48,91 ).
Entre os ingressantes, as maiores notas ocorreram também nas Regiões Nordeste ( 42,59 ) e Sul ( 45,18 ), todos acima da média nacional ( 41,72 ).
Novamente, para a análise segundo a Categoria Administrativa, todas as entidades \textbf{estão} representadas.
Observa -se comportamento semelhante àquele da Formação Geral, ou seja, as médias mais elevadas, tanto entre os concluintes como entre os ingressantes, foram encontradas em instituições de origem estadual ( 55,54 entre os concluintes e 51,80 entre os ingressantes ).
Os estudantes das instituições privadas apresentam uma maior homogeneidade.
Quanto à Organização Acadêmica, o comentário é o mesmo já feito para a nota bruta e referente à Formação Geral : as Universidades, Centros Universitários, Faculdades e Institutos Federais de Educação, Ciência e Tecnologia \textbf{estão} representados neste curso.
Considerando -se o tipo de Organização Acadêmica, entre os concluintes, Universidades apresentaram os maiores valores, seguidas das Faculdades e dos Centros Universitários.
E entre os ingressantes, Institutos Federais de Educação, Ciência e Tecnologia apresentaram os maiores valores, seguidos das Universidades, das Faculdades e, por fim, dos Centros Universitários. 3.2 Análise das Questões Objetivas 3.2.1 Formação Geral A Tabela 3.4 apresenta as estatísticas básicas relativas às oito questões objetivas da Formação Geral dos estudantes.
Adiantamos que a primeira questão foi desconsiderada pelo critério da correlação ponto-bisserial.
Como pode ser observado na Tabela 3.4, os estudantes concluintes obtiveram desempenho superior, com nota média de 51,29, em relação aos estudantes ingressantes cuja nota média foi 48,00.
A variabilidade entre os concluintes foi um pouco menor do que entre os ingressantes, como indicaram os desvios-padrão de 21,0 e 21,3, respectivamente.
A amplitude total das notas foi a mesma para ambos os grupos : nota mínima de 0,00 e máxima de 85,7.
A mediana dos concluintes foi 57,1 e a dos ingressantes, 42,8.
A assimetria, medida pelo coeficiente de assimetria ( skewness ) foi um pouco menor em módulo para os ingressantes ( -0,3 ) do que para os concluintes ( -0,4 ), ainda que ambos fossem negativos indicando uma assimetria para a direita, fato reforçado pelos valores menores das médias dos dois grupos quando comparados à mediana correspondente.
A Tabela 3.5 apresenta a classificação das questões objetivas da Formação Geral, segundo o índice de facilidade ( questões em vermelho foram desconsideradas pelo critério ponto-bisserial ).
Como as questões foram comuns a todas as áreas, a tabela apresenta a classificação considerando -se os concluintes de todas as áreas e aqueles do curso de COMUNICAÇÃO SOCIAL – PUBLICIDADE E PROPAGANDA.
De acordo com os índices obtidos, as questões objetivas da prova foram assim avaliadas para a área : das oito questões, duas tiveram entre 60 \%, inclusive, e 85 \% de acertos e foram classificadas como fáceis ( somente uma para o total de concluintes ) ; outras três foram consideradas medianas, alcançando entre 40 \%, inclusive, e 60 \% de acertos ( três para o total de concluintes ) ; duas obtiveram entre 15 \%, inclusive, e 40 \% de acertos, sendo consideradas difíceis ( três para o total de concluintes ).
Por fim, uma foi classificada como muito difícil tendo obtido percentual menor do que 15 \% de acertos ( uma para o total de concluintes ), coincidentemente a questão desconsiderada pelo critério ponto-bisserial.
Nenhuma questão foi classificada na categoria Muito fácil.
A Tabela 3.6 descreve os resultados das análises das questões objetivas relativas à Formação Geral, segundo o poder de discriminação, utilizando -se para tal o índice de discriminação ponto-bisserial.
Como as questões foram comuns a todas as áreas, a tabela apresenta a classificação considerando -se os concluintes de todas as áreas e aqueles do curso de COMUNICAÇÃO SOCIAL – PUBLICIDADE E PROPAGANDA.
A desconsideração da questão segundo este critério considerou o universo de todos os concluintes.
Duas das oito questões apresentaram índices iguais ou acima de 0,40 ( nenhuma para a população de concluintes como um todo ), sendo, portanto, classificadas como muito boas nesta característica, para este grupo de estudantes, uma foi classificada como boa com índice entre 0,30, inclusive, e 0,40 ( três para os concluintes como um todo ) e quatro foram classificadas como médias, com índice entre 0,20, inclusive, e 0,30 ( quatro para os concluintes como um todo ).
Uma das questões foi classificada como fraca tanto para a população de concluintes como um todo como para a de concluintes de COMUNICAÇÃO SOCIAL – PUBLICIDADE E PROPAGANDA e foi desconsiderada pelo critério. alcançando entre 40 \%, inclusive, e 60 \% de acertos ( três para o total de concluintes ) ; duas obtiveram entre 15 \%, inclusive, e 40 \% de acertos, sendo consideradas difíceis ( três para o total de concluintes ).
Por fim, uma foi classificada como muito difícil tento obtido percentual menor do que 15 \% de acertos ( uma para o total de concluintes ), coincidentemente a questão desconsiderada pelo critério ponto-bisserial.
Nenhuma questão foi classificada na categoria Muito fácil.
A Tabela 3.6 descreve os resultados das análises das questões objetivas relativas à Formação Geral, segundo o poder de discriminação, utilizando -se para tal o índice de discriminação ponto-bisserial.
Como as questões foram comuns a todas as áreas, a tabela apresenta a classificação considerando -se os concluintes de todas as áreas e aqueles do curso de COMUNICAÇÃO SOCIAL – PUBLICIDADE E PROPAGANDA.
A desconsideração da questão segundo este critério considerou o universo de todos os concluintes.
Duas das oito questões apresentaram índices iguais ou acima de 0,40 ( nenhuma para a população de concluintes como um todo ), sendo, portanto, classificadas como muito boas nesta característica, para este grupo de estudantes, uma foi classificada como boa com índice entre 0,30, inclusive, e 0,40 ( três para os concluintes como um todo ) e quatro foram classificadas como médias, com índice entre 0,20, inclusive, e 0,30 ( quatro para os concluintes como um todo ).
Uma das questões foi classificada como fraca tanto para a população de concluintes como um todo como para a de concluintes de COMUNICAÇÃO SOCIAL – PUBLICIDADE E PROPAGANDA e foi desconsiderada pelo critério.
Na sequência, no Quadro 3.1, \textbf{estão} especificadas, por questão, as habilidades aferidas pelas questões de Formação Geral da prova da COMUNICAÇÃO SOCIAL – PUBLICIDADE E PROPAGANDA, bem como os resultados da análise relativa aos índices de facilidade e de discriminação de \textbf{cada} questão.
O índice de facilidade para a área de COMUNICAÇÃO SOCIAL – PUBLICIDADE E PROPAGANDA teve uma variação considerável, oscilando entre 0,09 até 0,80 ( na população de concluintes o intervalo foi de 0,08 a 0,75 ) e o de discriminação ficou entre 0,02 e 0,48 ( na população de concluintes o intervalo foi de 0,02 a 0,39 ).
A questão 8 aferiu as habilidades de “ Ler e interpretar textos ”, “ Analisar e criticar informações ”, “ Extrair conclusões por indução e/ou dedução ”, “ Estabelecer relações, comparações e contrastes em diferentes situações ”, “ Detectar contradições ”, “ Fazer escolhas valorativas avaliando consequências ”, “ Questionar a realidade ”, e “ Argumentar coerentemente ” foi a que apresentou, para a população de concluintes de COMUNICAÇÃO SOCIAL – PUBLICIDADE E PROPAGANDA, maior poder discriminatório, ambas com índice de 0,48.
A questão mais Fácil, com 80 \% de acertos, foi também a de número 8.
De outra parte, a questão com menor índice de facilidade e de discriminação em Formação Geral foi a questão 1, que foi desconsiderada, e que deveria aferir as habilidades de “ Analisar informações ” ; “ Extrair conclusões ” e “ Propor soluções para situações-problema ”.
O Gráfico 3.7 mostra, a título de ilustração, o comportamento das respostas à questão 7 de Formação Geral para os estudantes concluintes da área de COMUNICAÇÃO SOCIAL – PUBLICIDADE E PROPAGANDA como função da nota geral na prova⁴.
Trata -se de uma questão considerada de dificuldade fácil para os concluintes de COMUNICAÇÃO SOCIAL – PUBLICIDADE E PROPAGANDA de dificuldade média para o conjunto de concluintes de todas as áreas.
O poder de discriminação foi muito bom para os estudantes da área e bom para o conjunto de todos os concluintes.
Neste Gráfico, \textbf{cada} uma das cinco curvas, representou o percentual de respostas em determinada alternativa da questão, em função do escore total dos estudantes concluintes da área de COMUNICAÇÃO SOCIAL – PUBLICIDADE E PROPAGANDA.
A curva azul marinho corresponde à alternativa A, a correta para esta questão.
Assim, observa -se que entre os estudantes com notas mais baixas – desempenho geral mais fraco na prova como um todo –, não há uma preponderância na escolha de nenhuma das alternativas, ainda que a alternativa correta seja já a mais frequente, seguida de alternativas incorretas, especialmente a alternativa B – curva verde claro e da alternativa C – curva em vermelho.
A proporção de estudantes que selecionaram a resposta correta A, aumentou com o desempenho deste grupo de estudantes chegando ao acerto de 100 \% dos estudantes com notas pelo menos 87,5.
Esta análise permitiu verificar como a questão discriminou entre os grupos de desempenho.
Os Gráficos relativos às demais questões de Formação Geral constam do Anexo I. ⁴ As notas da prova foram grupadas em 8 intervalos de amplitude 12,5, fechados à direita e abertos à esquerda e um ponto com a nota máxima. 3.2.2 Componente Específico A Tabela 3.7 apresenta as estatísticas básicas em relação às questões objetivas do Componente Específico.
Como pode ser observado, os estudantes concluintes obtiveram desempenho superior, com uma média de 52,78 pontos, em relação aos estudantes ingressantes com uma média de 45,54 pontos.
Os desvios-padrão indicaram que a variabilidade entre os concluintes foi maior ( 19,2 ) do que entre os ingressantes ( 17,5 ).
Os coeficientes de assimetria para os dois grupos foram negativos indicando uma inclinação para a direita.
A assimetria foi, porém, maior entre os concluintes.
A análise das questões objetivas do Componente Específico em COMUNICAÇÃO SOCIAL – PUBLICIDADE E PROPAGANDA, em relação ao grau de facilidade da prova ( ver Tabela 3.8 ), mostrou que, das 27 questões, nenhuma se encontrou na categoria ' muito difícil ' ( percentual de acertos menor do que 15 \% ) e nem na categoria ' muito fácil ' ( acima de 85 \% de acertos ), oito tiveram classificação de nível ' médio ' ( entre 41 \% e 60 \% de acertos ), quatro \textbf{estavam} na categoria ' fácil ' ( entre 61 \% e 85 \% de acertos ) e onze tiveram a classificação ' difícil ' ( entre 15 \% e 40 \% de acertos ).
Foram desconsideradas pela Comissão quatro questões.
Das 27 questões objetivas relativas ao Componente Específico ( ver Tabela 3.9 ), sete obtiveram o índice de discriminação muito bom, oito obtiveram o índice de discriminação bom e uma obteve nível de discriminação médio.
O índice fraco que eliminaria a questão foi alcançado por sete questões.
Foram desconsideradas pela comissão quatro questões.
Na sequência, no Quadro 3.2 \textbf{estão} especificadas, por questão, as habilidades aferidas pelas questões do Componente Específico da prova, bem como os resultados da análise relativa aos índices de facilidade e de discriminação de \textbf{cada} questão.
O Gráfico 3.8 mostra, a título de ilustração, o comportamento da questão 11 do Componente Específico para os estudantes concluintes da área de COMUNICAÇÃO SOCIAL – PUBLICIDADE E PROPAGANDA como função da nota geral na prova5.
Trata -se de uma questão considerada média para os concluintes de COMUNICAÇÃO SOCIAL – PUBLICIDADE E PROPAGANDA.
O poder de discriminação foi muito bom para os estudantes da área.
Neste Gráfico, \textbf{cada} uma das cinco curvas, representa o percentual de respostas em determinada alternativa da questão, em função da nota total dos estudantes concluintes de COMUNICAÇÃO SOCIAL – PUBLICIDADE E PROPAGANDA.
A curva verde claro correspondeu à alternativa B, a correta para esta questão.
Assim, observou -se que entre os estudantes com escores mais baixos – desempenho geral mais fraco na prova como um todo –, a tendência foi a escolha de uma alternativa incorreta, a alternativa A – curva azul marinho.
A proporção de estudantes que selecionaram a resposta correta B aumentou com o desempenho deste grupo de estudantes chegando ao acerto de 100 \% dos estudantes com notas pelo menos 87,5.
Esta análise permitiu verificar como a questão discriminou entre os grupos de desempenho.
Os Gráficos relativos às demais questões do Componente Específico constam do Anexo I. 3.3 Análise das Questões Discursivas 3.3.1 Formação Geral A Tabela 3.10 apresenta as estatísticas básicas da prova por grupo de estudantes.
A análise dos resultados de desempenho dos estudantes nas questões discursivas que são apresentados na tabela 3.10 e no gráfico 3.9 mostrou que as notas foram mais baixas no conjunto dessas questões que no das questões objetivas.
O total de estudantes inscritos para a prova foi de 36585.
Destes, 74,9 \% compareceram, sendo que a presença foi maior entre concluintes ( 85,2 \% ) do que entre ingressantes ( 66,4 \% ).
A média geral da prova foi de 25,42, sendo que os ingressantes obtiveram média mais baixa, de 24,29, que os concluintes, cuja média foi 26,48.
O desvio-padrão geral foi de 33,4, sendo o do grupo de ingressantes menor ( 32,6 ), em relação ao grupo dos concluintes ( 34,2 ), aqueles possuíram uma distribuição mais homogênea nas notas.
A nota máxima foi 100,0, atingida por concluintes e ingressantes.
Os coeficientes de assimetria para os dois grupos considerados, o de concluintes e o de ingressantes, bem como para a população como um todo foram todos positivos, indicando uma maior concentração à esquerda da distribuição.
Para todos os grupos, mais de metade dos estudantes presentes obtiveram nota zero nas questões discursivas da Formação Geral.
Das 16705 provas com nota zero, 602 ( 3,6 \% ) continham um protesto, 5289 ( 31,7 \% ) \textbf{estavam} em branco e 10814 ( 64,7 \% ) foram respondidas, porém de forma inadequada.
Entre os concluintes, a maioria ( 65,7 \% ) dentre aqueles com nota zero \textbf{estava} com as provas respondidas de forma inadequada.
Entre os ingressantes, 63,7 \% das notas nulas decorreram de correção.
Uma comparação entre o desempenho de ingressantes e concluintes pode ser realizada a partir dos dados apresentados no Gráfico 3.9.
As notas dos concluintes, de um modo geral, foram ligeiramente mais elevadas do que as dos ingressantes, cuja distribuição de notas \textbf{está} a esquerda dos primeiros.
A ideia de melhor desempenho foi reforçada pela proporção de notas nulas : 57,5 \% entre os concluintes e 59,4 \% entre os ingressantes.
Na sequência, discorrer -se-á, especificamente, sobre uma das questões discursivas de Formação Geral, a número 10 já que a número 9 foi desconsiderada.
Em especial, serão analisados os conteúdos englobados na questão e os resultados obtidos pelos estudantes.
Ao final, são apresentados os comentários da banca corretora a respeito do observado nas respostas dos estudantes, suas impressões e conclusões.
Cumpre esclarecer que, tendo em vista que as questões discursivas de Formação Geral foram padronizadas, ou seja, constam de todas as provas, os comentários da banca foram os mesmos para todas as áreas acadêmicas, eis que direcionados a todos os estudantes que participaram do ENADE/2009.
A seguir, serão analisados os desempenhos de ingressantes e concluintes da área de COMUNICAÇÃO SOCIAL – PUBLICIDADE E PROPAGANDA na questão discursiva de Formação Geral, não anulada, do ENADE/2009, comparando os resultados obtidos com a habilidade exigida na questão. 3.3.1.1 Análise da Questão Discursiva 9 de Formação Geral DESCONSIDERADA 3.3.1.2 Análise da Questão Discursiva 10 de Formação Geral A Questão Discursiva nº. 10 buscou avaliar as habilidades : • Ler e interpretar textos ; • Analisar e criticar informações ; • Extrair conclusões por indução e/ou dedução ; • Estabelecer relações, comparações e contrastes em diferentes situações ; • Detectar contradições ; • Fazer escolhas valorativas avaliando consequências ; • Questionar a realidade ; • Argumentar coerentemente.
O Quadro 3.3 apresenta as habilidades/competências, temas predominantes e estatísticas básicas da questão 10 de Formação Geral por grupo de estudantes.
Cumpre notar que como a questão 9 foi desconsiderada a nota referente às questões discursivas da Formação Geral é idêntica à nota da questão 10.
A média geral desta questão foi de 25,42, sendo que os ingressantes obtiveram média mais baixa, de 24,29, que os concluintes, cuja média foi 26,48.
O desvio-padrão geral foi de 33,4, sendo o do grupo de ingressantes menor ( 32,6 ), em relação ao grupo dos concluintes ( 34,2 ), aqueles possuíram uma distribuição mais homogênea nas notas.
A nota máxima foi 100,0, atingida por concluintes e ingressantes.
Os coeficientes de assimetria para os dois grupos considerados, o de concluintes e o de ingressantes, bem como para a população como um todo foram positivos, indicando uma maior concentração à esquerda da distribuição.
Para todos os grupos, cerca de metade dos estudantes presentes obtiveram nota zero nesta questões discursivas de Formação Geral.
Das 16705 provas com nota zero, 602 ( 3,6 \% ) continham um protesto, 5289 ( 31,7 \% ) \textbf{estavam} em branco e 10814 ( 64,7 \% ) foram respondidas, porém de forma inadequada.
Entre os concluintes, a maioria ( 65,7 \% ) dentre aqueles com nota zero \textbf{estavam} com as provas respondidas de forma inadequada.
Entre os ingressantes, 63,7 \% das notas nulas ( zero ) decorreram do desempenho do próprio estudante verificadas na correção.
Uma comparação entre o desempenho de ingressantes e concluintes para esta questão foi exatamente igual à já realizada para a nota discursiva de Formação Geral a partir dos dados apresentados no Gráfico 3.9. 3.3.1.3 Comentários sobre as respostas à Questão Discursiva 10 Apresentação A questão 10 de Formação Geral apresentou um questionamento capaz de analisar a habilidade do estudante em reconhecer o que é um direito social, optar por um direito social vigente que possa diminuir as desigualdades de renda familiar no Brasil e discutir os argumentos em favor do direito social escolhido como suporte à resposta.
Para análise das respostas apresentadas pelos estudantes, houve uma leitura inicial bastante criteriosa.
Foram considerados para efeito de pontuação, os seguintes argumentos : • Acesso à educação pública, gratuita e de qualidade, garantindo ao cidadão melhores condições de vida ; • Assistência à saúde, em seu contexto mais amplo, garantindo os cuidados essenciais de saúde, baseados em técnicas apropriadas ao alcance de todas as pessoas, contando com a participação da população ; • Permanência do estudante na escola, em todos os níveis escolares ( da educação infantil à educação superior, possibilitando a qualificação profissional ), o que garante melhores condições de trabalho e remuneração.
Considerou -se a qualificação e o desempenho docente, desenvolvendo nos estudantes melhores expectativas de aprendizagem e interesse pelo conteúdo da escola.
Os professores desenvolvendo metodologias adequadas e conhecimento pleno quanto a questões de aprendizagem ; • Condições dignas de trabalho, com remuneração que garanta qualidade de vida.
Foram consideradas as condições de trabalho, a preparação profissional para o trabalho e os cursos de aperfeiçoamento, oportunizando momentos de aprendizagem e desenvolvimento, a partir de situações desencadeadoras ; • Ser proprietário do imóvel em que reside, reduzindo os gastos com aluguel, garantindo melhores condições de vida, equilíbrio financeiro e construção do bem-estar familiar.
Desempenho dos estudantes na questão É importante salientar que havia respostas que desenvolveram adequadamente o questionamento apresentado, com fundamentação baseada em argumentos coerentes.
Contudo, em um número considerável de respostas, havia apenas um argumento e não dois, conforme fora solicitado.
Muitos estudantes informaram sobre problemas da educação, de um modo geral.
Algumas observações mencionadas pelos corretores durante o processo de correção foram : • há preocupação com uma educação de qualidade, com sugestões de cursos profissionalizantes e métodos preparatórios de aulas. • a educação é revelada como de baixa qualidade, exaltando -se treinamentos e cursos de atualização para a classe docente. • a saúde como essência para a vida.
O SUS é colocado como parte integrante de direito social, garantindo às pessoas fácil acesso. • oportunidade de emprego como forma de ascensão social, direito, cidadania. • os políticos brasileiros como causadores do despreparo educacional.
As respostas analisadas apresentaram norma culta, coesão e coerência razoáveis, que garantiram a compreensão da mensagem desenvolvida pelos estudantes.
Os erros ortográficos foram poucos e não comprometeram a discussão sobre o assunto e sobre os pontos de vista abordados.
Quanto à coerência, houve um desvio quanto ao objetivo da questão, já que alguns estudantes, através dos seus argumentos retrataram reivindicações, reclamações e situações de problemas enfrentados pelos brasileiros, citando golpes políticos.
Porém, mesmo quando se tratava desta avaliação, os estudantes foram autênticos através das explicações.
Houve erros quanto à concordância verbal e nominal.
Houve também, dificuldade em utilizar os conectores corretamente.
Alguns estudantes, por falta de ideias, apresentaram em suas respostas redundância e prolixidade.
O desempenho foi considerado satisfatório, pois a maioria dos estudantes soube responder ao questionamento.
Houve domínio do conteúdo, pois as ideias, os argumentos e a originalidade corresponderam ao esperado. 3.3.2 Componente Específico Nas questões discursivas do Componente Específico – questões de números 38, 39 e 40 –, a diferença de desempenho entre os grupos não foi muito acentuada, inclusive menor do que para as questões discursivas de Formação Geral.
A Tabela 3.11 contém as estatísticas da média destas questões.
Observa -se que no grupo dos concluintes, a nota média alcançada foi 27,34, enquanto que no grupo dos ingressantes foi 20,37, ambas menores do que as obtidas nas questões discursivas de Formação Geral, que foram 23,97 para os ingressantes e 26,5 para os concluintes ( Tabela 3.10 ).
As médias foram também menores do que as obtidas para as questões objetivas do Componente Específico ( Tabela 3.7 ) : respectivamente 45,5 e 52,8 para ingressantes e concluintes.
O desempenho mais baixo nas questões discursivas refere -se não só à maior complexidade desse tipo de questão 38 –, a diferença de desempenho entre os grupos não foi muito acentuada, inclusive menor do que para as questões discursivas de Formação Geral.
A Tabela 3.11 contém as estatísticas da média destas questões.
Observa -se que no grupo dos concluintes, a nota média alcançada foi 27,34, enquanto que no grupo dos ingressantes foi 20,37, ambas menores do que as obtidas nas questões discursivas de Formação Geral, que foram 23,97 para os ingressantes e 26,5 para os concluintes ( Tabela 3.10 ).
As médias foram também menores do que as obtidas para as questões objetivas do Componente Específico ( Tabela 3.7 ) : respectivamente 45,5 e 52,8 para ingressantes e concluintes.
O desempenho mais baixo nas questões discursivas refere -se não só à maior complexidade desse tipo de questão, mas também a uma dificuldade, por parte dos estudantes de COMUNICAÇÃO SOCIAL – PUBLICIDADE E PROPAGANDA, de discorrer sobre temas específicos da área.
A análise de \textbf{cada} uma destas questões será feita nas subseções seguintes.
Os valores positivos para a medida de assimetria indicaram uma distribuição com cauda para a direita.
Dos 27418 presentes, 28,7 \% obtiveram nota zero.
A tabela 3.12 apresenta um \textbf{detalhamento} da situação destes estudantes.
A maioria ( 48,2 \% ) não fez nenhuma das 3 questões.
Esta proporção foi menor entre os concluintes ( 47,6 \% ) do que entre os ingressantes ( 48,7 \% ).
A combinação de questões em branco e preenchidas correspondeu a 96,7 \% das provas.
Menos de 1 \% no total das provas correspondeu a 3 questões com protesto.
O Gráfico 3.10 apresenta a distribuição das notas nas questões discursivas do Componente Específico por grupo de estudantes.
A moda se situou no primeiro intervalo, aquele com valores entre 0 e 10 incluindo os extremos do intervalo.
A impressão de assimetria informada pelos coeficientes de assimetria positivos ( Tabela 3.11 ) foi reforçada no gráfico pela concentração no primeiro intervalo. 3.3.2.1 Análise da Questão Discursiva 38 do Componente Específico A questão 38, cujas habilidades aferidas encontram -se descritas no Quadro 3.4, que ainda contém outras informações, seguido da Tabela 3.13 e do Gráfico 3.11, foi aquela em que os grupos de estudantes observados obtiveram o melhor desempenho entre as questões discursivas do Componente Específico.
A nota média do grupo dos ingressantes foi de 20,90, e a dos concluintes 32,89.
O percentual de notas zero obtidas, o mais baixo das questões discursivas do Componente Específico, nos dois grupos foi de 39,7 \%, maior entre os ingressantes ( 45,1 \% ) do que entre os concluintes ( 34,6 \% ).
Quase 70 \% dos concluintes e dos ingressantes que tiraram zero deixaram a questão em branco ( ver Tabela 3.13 ).
Uma pequena minoria ( 2,9 \% ) utilizou a questão para protestar.
Apesar do baixo desempenho, em ambos os grupos a nota máxima de 100,0 pontos foi alcançada.
Os coeficientes de assimetria positivos indicam uma inclinação para a esquerda e cauda para a direita.
O Gráfico 3.11 apresenta a distribuição das notas na questão 38 discursiva do Componente Específico por grupo de estudantes.
A moda se situou no primeiro intervalo, aquele com valores entre 0 e 10 incluindo os extremos do intervalo.
A impressão de assimetria informada pelos coeficientes de assimetria positivos ( Quadro 3.4 ) foi reforçada no gráfico pela concentração no primeiro intervalo. 3.3.2.2 Comentários sobre as respostas à Questão Discursiva 38 O objetivo central da questão foi \textbf{apurar} o conhecimento dos estudantes sobre as novas possibilidades de uso dos meios eletrônicos para a publicidade.
O próprio enunciado já procurava direcionar a resposta para a solução correta ao expor dados sobre o problemático cenário econômico mundial pós-crise, a redução das verbas publicitárias decorrentes dele e o baixo custo de veiculação de publicidade em meio digital.
Quando os estudantes demonstraram ter compreendido o comando e o objetivo do enunciado, em geral, obtiveram um desempenho de razoável para bom, mas o percentual de respostas corretas foi bem mais baixo que o esperado, uma vez que o próprio enunciado direcionava a resposta para a solução correta : um planejamento voltado aos meios digitais e seus recursos.
Houve um bom percentual de respostas que apresentaram soluções como a articulação entre as opções digitais e as tradicionais de mídia, com predomínio das primeiras, exatamente como seria o desejável.
As respostas, muitas vezes, não se enquadraram às limitações impostas, talvez por não serem muito frequentes nesta área, ações propostas exclusivamente às mídias digitais, já que estas são tradicionalmente as que recebem menos recursos.
Também foi frequente nas respostas, o erro de querer estabelecer condições financeiras para a venda dos imóveis ( prazo de financiamento, valores de juros ou mesmo inexistência deles, isenção do pagamento de parcelas como promoção, etc. ) ou determinar características arquitetônicas específicas ( acabamentos luxuosos, recursos modernos de respeito ao meio ambiente, existência de áreas de lazer, etc. ), apesar de nada disso ser citado no enunciado, ou solicitado, e de todos esses aspectos fugirem totalmente às atribuições de um profissional de publicidade. 3.3.2.2 Comentários sobre as respostas à Questão Discursiva 38 O objetivo central da questão foi \textbf{apurar} o conhecimento dos estudantes sobre as novas possibilidades de uso dos meios eletrônicos para a publicidade.
O próprio enunciado já procurava direcionar a resposta para a solução correta ao expor dados sobre o problemático cenário econômico mundial pós-crise, a redução das verbas publicitárias decorrentes dele e o baixo custo de veiculação de publicidade em meio digital.
Quando os estudantes demonstraram ter compreendido o comando e o objetivo do enunciado, em geral, obtiveram um desempenho de razoável para bom, mas o percentual de respostas corretas foi bem mais baixo que o esperado, uma vez que o próprio enunciado direcionava a resposta para a solução correta : um planejamento voltado aos meios digitais e seus recursos.
Houve um bom percentual de respostas que apresentaram soluções como a articulação entre as opções digitais e as tradicionais de mídia, com predomínio das primeiras, exatamente como seria o desejável.
As respostas, muitas vezes, não se enquadraram às limitações impostas, talvez por não serem muito frequentes nesta área, ações propostas exclusivamente às mídias digitais, já que estas são tradicionalmente as que recebem menos recursos.
Também foi frequente nas respostas, o erro de querer estabelecer condições financeiras para a venda dos imóveis ( prazo de financiamento, valores de juros ou mesmo inexistência deles, isenção do pagamento de parcelas como promoção, etc. ) ou determinar características arquitetônicas específicas ( acabamentos luxuosos, recursos modernos de respeito ao meio ambiente, existência de áreas de lazer, etc. ), apesar de nada disso ser citado no enunciado, ou solicitado, e de todos esses aspectos fugirem totalmente às atribuições de um profissional de publicidade.
Concluindo, podemos afirmar que o enfoque foi de agrado dos estudantes e não foram registrados impropérios ou desabafos, quando da resposta a questão. 3.3.2.3 Análise da Questão Discursiva 39 do Componente Específico A questão 39, cujas habilidades aferidas encontram -se descritas no Quadro 3.5, que ainda contém outras informações, seguido da Tabela 3.14 e do Gráfico 3.12, foi aquela em que os grupos de estudantes observados obtiveram desempenho intermediário entre as questões discursivas do Componente Específico.
A nota média do grupo dos ingressantes foi de 20,84, e a dos concluintes 26,00.
O percentual de notas zero obtidas nos dois grupos foi de 50,5 \% entre os ingressantes e 48,1 \% entre os concluintes.
Das provas com notas nulas nesta questão, um pouco mais de 1/3 foi deixada em branco : 37,6 \% entre os ingressantes e 31,6 \% entre os concluintes ( ver Tabela 3.14 ).
Entre os com nota nula, uma pequena minoria ( 1,2 \% ) utilizou a questão para protestar.
Apesar do baixo desempenho, em ambos os grupos a nota máxima de 100,0 pontos foi alcançada.
Os coeficientes de assimetria positivos indicaram uma inclinação para a esquerda e cauda para a direita.
O Gráfico 3.12 apresenta a distribuição das notas na questão 39 discursiva do Componente Específico por grupo de estudantes.
A moda se situou no primeiro intervalo, aquele com valores entre 0 e 10 incluindo os extremos do intervalo.
A impressão de assimetria informada pelos coeficientes de assimetria positivos ( Quadro 3.5 ) foi reforçada no gráfico pela concentração no primeiro intervalo.
A distribuição apresentou um máximo local no terceiro intervalo, de notas baixas, indicando um baixo aproveitamento dos estudantes. 3.3.2.4 Comentários sobre as respostas à Questão Discursiva 39 A questão pode ser considerada fácil, exigindo dos estudantes apenas uma análise baseada em suas próprias experiências e conhecimentos práticos, sem exigência acentuada de conhecimentos técnicos ou teóricos da área.
Bastou -se perceber que os ícones apresentados representaram outra geração, distinta e almejada, portanto, não se enquadraram, e que a distribuição das mídias flyers e folders em locais ( boites e bares ) onde as pessoas não \textbf{estão} propícias a recepção deste material era equivocada.
O principal fator para a correção da resposta e um bom desempenho era a percepção da inadequação de atores famosos, principalmente nas décadas de 1940, 1950 e 1960 com um perfil de consumidor com faixa etária entre 25 e 35 anos.
Quando esse dado era claramente percebido pelo estudante, o desempenho era, no mínimo, satisfatório com alto índice de notas elevadas entre esse grupo.
Houve um percentual razoável de respostas focadas exclusivamente na correção do dado do planejamento de mídia que serviu de erro secundário.
O erro causado pela incompreensão do correto significado da expressão “ esporte fino ”, confundido com as vestimentas utilizadas para a prática esportiva de modalidades, como o tênis e o golfe ( citadas frequentemente ) foi bastante corriqueiro, assim como a opção por esportistas e/ou outros atores sociais que não são considerados modelos de elegância.
Acredita -se que o desconhecimento do significado dessa expressão demonstrou a dificuldade de vivência social de boa parcela dos estudantes.
Outro erro, mais grave, foi decorrente do total desconhecimento das condições de vida da população pertencente à classificação socioeconômica descrita no enunciado e ao preconceito.
Diversas respostas apontaram como erro a inclusão das classes sociais C+ e B- entre as que comporiam o público alvo da empresa, muitas vezes essa afirmação era justificada por expressões como “ pessoas das classes sociais C+ e B- não podem comprar roupas de boa qualidade ”. 3.3.2.5 Análise da Questão Discursiva 40 do Componente Específico A questão 40, cujas habilidades aferidas encontram -se descritas no Quadro 3.6, que ainda contém outras informações, seguido da Tabela 3.15 e do Gráfico 3.13, foi aquela em que os grupos de estudantes observados obtiveram o pior desempenho entre as questões discursivas do Componente Específico.
A nota média do grupo dos ingressantes foi de 19,45, e a dos concluintes 23,19.
Nesta questão ocorre também o maior percentual de notas nulas : 55,7 \% entre os ingressantes e 52,5 \% entre os concluintes.
Aqui a maioria das provas com nota zero decorreu de provas com respostas desenvolvidas inadequadamente : 59,9 \% entre os ingressantes e 64,2 \% entre os concluintes ( ver Tabela 3.15 ).
Uma pequena minoria ( 1,1 \% ), dentre aqueles com nota nula, utilizou a questão para protestar.
Em ambos os grupos a nota máxima de 100,0 pontos foi alcançada.
Os coeficientes de assimetria positivos indicaram uma inclinação para à esquerda.
O Gráfico 3.13 apresenta a distribuição das notas na questão 40 discursiva do Componente Específico por grupo de estudantes.
A moda se situou no primeiro intervalo, aquele com valores entre 0 e 10 incluindo os extremos do intervalo.
A impressão de assimetria informada pelos coeficientes de assimetria positivos ( Quadro 3.6 ) foi reforçada no gráfico pela concentração no primeiro intervalo. 3.3.2.6 Comentários sobre as respostas à Questão Discursiva 40 O objetivo central da questão era \textbf{apurar} o conhecimento dos estudantes sobre os processos e as ferramentas de comunicação integrada.
O enunciado exigiu do estudante uma boa capacidade de compreensão do processo de comunicação como um todo, direcionando a resposta para ações de reposicionamento de marca que deveriam ser praticadas, para que a comunicação da empresa-problema fosse integralmente reformulada, visando à conquista de um novo público alvo, uma vez que a razão alegada para a perda de público consumidor anterior era o seu envelhecimento.
Esta foi a questão com menor aproveitamento entre as demais.
Houve grande dificuldade dos estudantes de perceberem que o solicitado pelo enunciado não era a simples criação de peças de divulgação da marca.
Mesmo entre os estudantes que procuraram atender ao comando, ficou patente o pequeno grau de conhecimento do ponto em foco ( comunicação integrada ).
O erro mais frequente foi a apresentação de respostas, cujo conteúdo era a simples sugestão da produção de algumas peças publicitárias.
Também foi comum a sugestão de investimento maior na reconquista do público consumidor perdido.
Uma vez que a razão alegada para perda desse público era seu envelhecimento e que a empresa-problema era uma fábrica de bebidas e que deveria ser presumível pelo estudante que tal investimento fosse inócuo, somente pela leitura atenta do enunciado, o que deveria gerar o consequente questionamento das razões que levariam um público envelhecido a abandonar o hábito do consumo de bebidas alcoólicas.
Outro equívoco frequente foi a sugestão da feitura de uma pesquisa de mercado, dado já informado no enunciado como realizado.
Nesse caso, a desatenção ao enunciado provocou o erro, uma vez que, na maioria dessas respostas não ocorreu a mera repetição do enunciado, mas sim o desenvolvimento deste dado como sendo algo novo.
Além dos problemas já apontados, também foi comum a compreensão errônea de que a empresa-problema ( produtora de bebidas alcoólicas ) seria uma cervejaria, o que direcionou a resposta para as soluções publicitárias mais visíveis neste segmento do mercado publicitário contemporâneo.
A escrita, de modo geral, apresentou erro de ortografia, o que foi evidenciado nos argumentos apresentados.

% matched lemmas: apurar, cada, detalhamento, estar
\end{textsample}
